% Baccalauréat malgache, série D, 2018 : énoncé.
\TCDBacAttribution{D}{2018}
\providecommand{\Plan}{(\mathscr{P})}
\providecommand{\Ouv}{(O ; \vec{u}, \vec{v})}
\providecommand{\cm}{~\mathrm{cm}}
\providecommand{\Courbe}{(\mathscr{C})}
\providecommand{\CourbeP}{(\mathscr{C}')}
\providecommand{\D}{(\mathscr{D})}
\providecommand{\Droite}{(\mathscr{D})}
\providecommand{\dx}{~\mathrm{d}x}
\providecommand{\Angl}[2]{\left( \overrightarrow{#1} ; \overrightarrow{#2} \right)}
%%%%%%%%%%%%%%%%%%%%%%%%%%%%%%%%%%%%%%%%%%%%%%%%%%%%%%%%%%%%%%%%%%%%%%%%%%%%%%%%%%%%%%%%%%%%%%%
%%%%%%%%%%%%%%%%%%%                   %%%%%%%%%%%%%%%%%%%                   %%%%%%%%%%%%%%%%%%%
%%%%%%%%%%%%%%%%%%%%%%%%%%%%%%%%%%%%%%%%%%%%%%%%%%%%%%%%%%%%%%%%%%%%%%%%%%%%%%%%%%%%%%%%%%%%%%%
\noindent \underline{\textsc{Exercice I}} (5 points) 
      Soit $P$ le polyn\^ome \`a variable complexe $z$ d\'efini par :
      \[
            P(z) = z^3 + (1 - 9i) z^2 - (24 + 6i) z - 14 + 18 i \;.
      \]
      \begin{enumerate}
            \item \begin{enumerate}
                        \item D\'eterminer le nombre complexe $z_0$ tel que : 
                              \[
                                    P(z) = (z - z_0) \; ( z^2 - 4 i z - 4 - 2i ) \;.
                              \]
                        \item En d\'eduire les solutions dans $\mathbb{C}$ de l'\'equation $P(z) = 0$.
                  \end{enumerate}
            \item Dans le plan complexe $\Plan$ muni d'un rep\`ere orthonorm\'e direct $\Ouv$ d'unit\'e $1 \cm$,
                  on consid\`ere les points $A$, $B$ et $C$ d'affixes respectives : 
                  \[
                        z_A = -1 + 5i \ ; \ z_B = 1 + 3i \text{ et } z_C = -1 + i \;.
                  \]
                  \begin{enumerate}
                        \item Placer les points $A$, $B$ et $C$.
                        \item Calculer les distances $AB$ et $BC$ et d\'eterminer une mesure de l'angle $\Angl{BA}{BC}$.
                              En d\'eduire la nature du triangle $ABC$.
                        \item On note $I$ le milieu du segment $[AC]$. 
                              D\'eterminer l'affixe $z_I$ du point $I$.
                        \item D\'eterminer puis construire l'ensemble $\Courbe$ des points $M$ d'affixe $z$ tels que 
                              $|z + 1 - 3i| = 2$.
                  \end{enumerate}
            \item Soit $S$ la similitude plane directe qui laisse invariant le point $A$ et transforme le point $C$ en $B$.
                  \begin{enumerate}
                        \item D\'eterminer l'expression complexe ainsi que les \'el\'ements caract\'eristiques de $S$.
                        \item D\'eterminer et construire l'ensemble $\CourbeP$ image de $\Courbe$ par $S$.
                  \end{enumerate}
      \end{enumerate}
%%%%%%%%%%%%%%%%%%%%%%%%%%%%%%%%%%%%%%%%%%%%%%%%%%%%%%%%%%%%%%%%%%%%%%%%%%%%%%%%%%%%%%%%%%%%%%%
%%%%%%%%%%%%%%%%%%%                   %%%%%%%%%%%%%%%%%%%                   %%%%%%%%%%%%%%%%%%%
%%%%%%%%%%%%%%%%%%%%%%%%%%%%%%%%%%%%%%%%%%%%%%%%%%%%%%%%%%%%%%%%%%%%%%%%%%%%%%%%%%%%%%%%%%%%%%%
\noindent \underline{\textsc{Exercice II}} (5 points) 
      \begin{enumerate}
            \item On lance une fois un d\'e t\'etra\'edrique \`a quatre faces num\'erot\'ees 1, 2, 3 et 4.
                  
                  On s'interesse au num\'ero de la face cach\'ee.
                  Pour $k \in \{ 1; 2 ; 3 ; 4 \}$, on note $P_k$ la probabilit\'e pour que 
                  le num\'ero de la face cach\'ee soit \'egale \`a $k$. \newline
                  Le d\'e est truqu\'e de telle sorte que les probabilit\'es $P_1$, $P_2$, $P_3$ et $P_4$ 
                  v\'erifient les conditions suivantes : 
                  \[
                        P_2 = \frac{1}{2} P_1 = \frac{1}{3} P_3 \text{ et } P_4 = 2 P_1 \;.
                  \]
                  D\'emontrer que $P_1 = \frac{1}{5}$.
                  En d\'eduire les probabilit\'es $P_2$, $P_3$ et $P_4$.
            \item On lance deux fois de suite ce m\^eme d\'e.
                  \begin{enumerate}
                        \item Calculer la probabilit\'e de l'\'ev\`enement : \newline
                              $A$ : << Le produit des deux num\'eros des deux faces cach\'ees est \'egal \`a 4 >>.
                        \item On d\'esigne par $a$ le num\'ero de la face cach\'ee au premier lancer et par $b$ 
                              le num\'ero du deuxi\`eme.
                              Soit $X$ la variable al\'eatoire \'egale au nombre $| b - a|$.
                              \begin{itemize}
                                    \item Donner la loi de probabilit\'e de $X$.
                                    \item Calculer l'esp\'erance math\'ematique $E(X)$ de $X$.
                              \end{itemize}
                  \end{enumerate}
            \item On lance quatre fois de suite et d'une mani\`ere ind\'ependante ce d\'e.
                  Soit $Y$ la variable al\'eatoire \'egale au nombre de fois o\`u la 
                  face n\textsuperscript{$\circ$}1 est cach\'ee.
                  \begin{enumerate}
                        \item Calculer la probabilit\'e de l'\'ev\`enement $(Y \geq 1)$.
                        \item Calculer $E(Y)$ et $V(Y)$.
                  \end{enumerate}
      \end{enumerate}
%%%%%%%%%%%%%%%%%%%%%%%%%%%%%%%%%%%%%%%%%%%%%%%%%%%%%%%%%%%%%%%%%%%%%%%%%%%%%%%%%%%%%%%%%%%%%%%
%%%%%%%%%%%%%%%%%%%                   %%%%%%%%%%%%%%%%%%%                   %%%%%%%%%%%%%%%%%%%
%%%%%%%%%%%%%%%%%%%%%%%%%%%%%%%%%%%%%%%%%%%%%%%%%%%%%%%%%%%%%%%%%%%%%%%%%%%%%%%%%%%%%%%%%%%%%%%
\noindent \underline{\textsc{Probleme}} \newline
      On consid\`ere la fonction num\'erique $f$ d\'efinie sur $]0 ; +\infty[$ par : 
      \[
            f(x) = - x + 3 + \frac{(-1+\ln x)}{x} \;.
      \]
      On note par $\Courbe$ sa courbe repr\'esentative dans un rep\`ere orthonorm\'e \Oij{} d'unit\'e $1\cm$.
      \begin{enumerate}
            \item Soit $g$ la fonction d\'efinie sur $]0 ; +\infty[$ par : $g(x) = x^2 - 2 + \ln x$.
                  \begin{enumerate}
                        \item Etudier le sens de variation de $g$ sur $] 0 ; +\infty[$.
                        \item Montrer que l'\'equation $g(x) = 0$ admet une solution r\'eelle unique $\alpha$ telle que
                              $\alpha \in ] 1,2\ ;\ 1,4[$.
                        \item En d\'eduire le signe de $g(x)$ suivant les valeurs de $x$.
                  \end{enumerate}
            \item \begin{enumerate}
                        \item Calculer $\displaystyle{\lim_{x\rightarrow 0^+} f(x)}$.
                              Interpr\'eter graphiquement le r\'esultat.
                        \item Calculer $\displaystyle{\lim_{x\rightarrow +\infty} f(x)}$.
                        \item Montrer que la droite $\Droite$ d'\'equation $y = -x + 3$ est asymptote oblique \`a $\Courbe$.
                        \item Etudier la position relative de $\Courbe$ par rapport \`a $\Droite$.
                  \end{enumerate}
            \item \begin{enumerate}
                        \item Montrer que, pour tout $x \in ]0 ; +\infty[$, 
                              $\displaystyle{f'(x) = \frac{-g(x)}{x^2} }$.
                        \item D\'emontrer que $\displaystyle{f(\alpha) = \frac{-2\alpha^2 + 3 \alpha + 1}{\alpha} }$.
                        \item Dresser le tableau de variation de la fonction $f$ sur $]0 ; +\infty[$.
                  \end{enumerate}
            \item \begin{enumerate}
                        \item D\'eterminer les coordonn\'ees du point $A$ de $\Courbe$ o\`u la tangente $(T)$ est
                              parall\`ele \`a $\Droite$.
                        \item Construire $(T)$, $\Droite$ et $\Courbe$.
                              (On prendra $\alpha \simeq 1,3$ pour la construction).
                  \end{enumerate}
            \item Calculer, en $\cm^2$, l'aire du domaine $\mathcal{A}$ du domaine plan d\'elimit\'e par la courbe $\Courbe$,
                  la droite $\Droite$ et les droites d'\'equations $x = 1$ et $x = e$.
      \end{enumerate}
      %
      \begin{center}
            \underline{On donne :} $e \simeq 2,71$ ; $\expo{2} \simeq 7,38$ ; $\frac{1}{\expo{2}} \simeq 0,13$.
      \end{center}











